\documentclass{ltjsbook}
\usepackage{docmute} 
\usepackage{../../myStyle} 

\begin{document}
\chapter{固有値と固有ベクトル}
\label{L07_Eigen}
\section{固有値と固有ベクトル}
今回は正方行列にみられるおもしろい特徴である，\keyterm{固有値}{こゆうち}{eigenvalue}と\keyterm{固有ベクトル}{こゆうべくとる}{eigenvector}についての話から始めます。
ある正方行列$\bm{A}$，列ベクトル$\bm{x}$，スカラー$\lambda$が次のような関係にあった時，$\lambda$を固有値，$\bm{x}$を固有ベクトルと言います。

\[ \bm{Ax} = \lambda\bm{x} \]

一見すると，$\bm{x}$が両辺に入っていますから，$\bm{A}$が$\lambda$に置き換わった等式のようです。しかし一方は行列で，他方はスカラーです。こんな奇妙なことが本当にあるのでしょうか？
具体的な数値例をみてみましょう。

\[
	\bm{A} = \begin{pmatrix}
		1 & 6 \\ 2 & 5
	\end{pmatrix}
	,\bm{x} = \begin{pmatrix}
		1 \\ 1
	\end{pmatrix}
\]

を例にします。この時次の関係が成り立ちます。

\[
	\bm{Ax} =
	\begin {pmatrix}
	1 & 6 \\
	2 & 5 \\
	\end{pmatrix}
	\begin{pmatrix}
		1 \\ 1
	\end{pmatrix}
	=
	\begin{pmatrix}
		7 \\ 7
	\end{pmatrix}
	=7
	\begin{pmatrix}
		1 \\ 1
	\end{pmatrix}
	=7 \bm{x}
\]

確かに成立する組み合わせがありますね。この行列$\bm{A}$に対して，7が固有値，$(1,1)$が固有ベクトルになっています。
また，実はこの行列$\bm{A}$については，-1も固有値であり，そのときの固有ベクトルは$(-3,1)$になります。

この固有値分解こそ，因子分析を元とする多変量解析の中心的な数学原理なのです。多変量解析の世界においては，分散共分散行列やそれを標準化した相関行列など，変数同士の関係を分析のスタートにおくのでした。これらの行列は正方行列ですから，その固有値や固有ベクトルを計算することで正方行列の特徴を別の視点から分解して考えられるようになります。



\subsection{固有値と固有ベクトルを求める}
では正方行列からどのようにして固有値と固有ベクトルを見つけ出したら良いのか，みていきましょう。
元の式を書き換えて次のような方程式を考えます。

\[ (\bm{A} - \lambda\bm{I})\bm{x} = \bm{0} \]

固有ベクトルは$\bm{x}=\bm{0}$すなわち全部ゼロであれば当然成り立ちますから(自明な解)，これは除外することにします($\bm{x}\neq \bm{0}$)。

さて行列の表現は連立方程式の解を求めることと同じなのでした。$\bm{A} - \lambda\bm{I}$を連立方程式の係数行列だと考えれば，それが\keytermL{逆行列}{ぎゃくぎょうれつ}を持つと左辺にそれをかけてしまえば全部ゼロの答えになってしまいますから，そうでない答えを求めるには，この係数行列が逆行列を持たないことが重要です。そして逆行列を持たないようにするというのはつまり，行列式がゼロであれば良いことになります。

$2 \times 2$の係数行列，$\bm{P}=\begin{pmatrix} a&b\\c&d \end{pmatrix}$の逆行列は$ad-bc$で求められるのでした。
$ad-bc$が0でなければ方程式は解けるのですが，今回の場合は解けると自明になってしまうので困ります。今回は$ad-bc=0$でなければならないのです。つまり一般的に書くと次のようになります。

\[ | \bm{A} - \lambda \bm{I} | = 0 \]

$\bm{A} = \begin{pmatrix} a & b \\ c & d \end{pmatrix} $とすると，この式は次のようになります。

\[   \begin{vmatrix} a-\lambda & b \\ c & d - \lambda \end{vmatrix} = 0 \]
\[ (a-\lambda)(d-\lambda) - bc = 0 \]

この方程式をとくに\keytermL{固有方程式}{こゆうほうていしき}といいますが，これを解いてやれば良いことになりますね。具体的に$\begin{pmatrix}1&6\\2&5\end{pmatrix}$の例で計算してみましょう。
\[ \begin{vmatrix}1-\lambda & 6 \\ 2 & 5-\lambda\end{vmatrix}=0\]
\[ (1-\lambda)(5-\lambda) - 12 = 0 \]
\[ \lambda^2 - 6\lambda -7 = 0 \]
\[ (\lambda-7)(\lambda+1)=0\]
ここから$\lambda=7,-1$が得られますね。

では固有ベクトルはどうなるでしょうか。固有値7の例で計算してみます。
\[ \begin{pmatrix} 1 & 6 \\ 2 & 5 \end{pmatrix}\begin{pmatrix} x\\y\end{pmatrix} = 7 \begin{pmatrix} x\\y\end{pmatrix} \]
\[ x + 6y = 7x \to 6x = 6y \to x=y \]
\[ 2x+5y = 7y \to 2x = xy \to x=y \]

あれっ？ なんじゃこりゃ？ と思った人もいるかもしれません。でもこれ，間違いではありません。実は固有ベクトルは大きさが定まっておらず，今回の例で言えば$x=y$つまり$x:y=1:1$の関係であればあらゆる値が成立してしまうのです。
固有ベクトルが$(1,1)$や$(2,2)$，$(100,100)$などであっても，$\bm{Ax}=\lambda \bm{x}$の関係に影響しませんから，計算機で算出する場合などは，ベクトルのノルムを1.0に規格化するという条件を追加することが一般的です。

\begin{reidai}{}{}
	行列$\begin{pmatrix}8&1\\4&5\end{pmatrix}$の固有値・固有ベクトルを求めてみましょう。
\end{reidai}
さてここでみたように，$2\times2$の方程式であれば手計算でも固有方程式を解くことはできるのですが，行列のサイズがどんどん大きくなると一般的に解けなくなって行くことは想像にかたくないと思います。実際我々は正方行列として，項目の情報が詰まった分散共分散行列とか相関行列を使いますから，それが2項目しかないなんてことはなくて，もっともっと大きなサイズになります。そうすると計算機を使って近似的に答えを求めて行くことになります。

\subsection{固有値と固有ベクトルの性質}
ではここで，$n$次正方行列$\bm{A}$の固有値と固有ベクトルに関する性質をみていきましょう。
\paragraph{固有値の重複度}

固有値は固有方程式から求められますが，固有方程式は$n$次方程式であり，解は$n$個あることになります。これは代数学の基本定理\footnote{「$p$次方程式は$p$個の根を持つ」ことが代数学の基本定理ですが，証明には複素数も含めた議論が必要ですのでここでは省略します。ごめんなさい！}が示すことでもあります。

ただし，同じ解が重複して出ること(重解)もあることに注意です。正確には「$n$次正方行列の固有ベクトルは，重複度を入れて$n$個ある」ということです。

\paragraph{固有ベクトルは一次独立}

行列$\bm{A}$の異なる固有値に対する固有ベクトルは，一次独立です。証明してみましょう。
$\bm{A}$の固有値$\lambda_1,\lambda_2,\cdots,\lambda_p$がそれぞれ異なり，対応する固有ベクトルが$\bm{x}_1,\bm{x}_2,\cdots,\bm{x}_p$だったとします。もし$\bm{x}_1,\bm{x}_2,\cdots,\bm{x}_p$が一次独立でないとすると，あるベクトル$\bm{x}_i$はほかの固有ベクトルの一次結合で表せることになりますね。たとえば$\bm{x}_1$が残りの一次独立な固有ベクトル$,\bm{x}_{2},\bm{x}_{3},\cdots,\bm{x}_{p}$ で表せたとします。つまり，
\[ \bm{x}_1 = c_1\bm{x}_{2} + c_2\bm{x}_{3} + \cdots + c_{p-1}\bm{x}_{p} \]
ということですね。これに行列$\bm{A}$を左からかけると，
\[ \bm{A}\bm{x}_1 = \bm{A}(c_1\bm{x}_{2} + c_2\bm{x}_{3} + \cdots + c_{p-1}\bm{x}_{p}) \]
より
\[ \lambda_1\bm{x}_1= c_1\bm{A}\bm{x}_{2} + c_2\bm{A}\bm{x}_{3} + \cdots + c_{p-1}\bm{A}\bm{x}_{p} \]
\[ = c_1\lambda_{2}\bm{x}_{2} + c_2\lambda_{3}\bm{x}_{3} + \cdots + c_{p-1}\lambda_{p}\bm{x}_{p} \]
となります。さて，一次結合の式に$\lambda_1$をかけてやると
\[ \lambda_1\bm{x}_1 = c_1\lambda_1\bm{x}_{2} + c_2\lambda_1\bm{x}_{3} + \cdots + c_{p-1}\lambda_1\bm{x}_{p} \]
ですから，両辺それぞれ引き算すれば，
\[ \bm{0} = c_1(\lambda_{2}-\lambda_1)\bm{x}_{2} + c_2(\lambda_{3}-\lambda_1)\bm{x}_{3} + \cdots +c_{p-1}(\lambda_{p}-\lambda_1)\bm{x}_{p} \]
となります。固有値はすべて異なるということでしたから，$c_1=c_2=\cdots=c_{p-1}=0$となります。さてそうすると，$\bm{x}_1=\bm{0}$となってしまいますが，固有ベクトルはゼロではないのでこれは矛盾です。したがって，$\bm{x}_1,\bm{x}_2,\cdots,\bm{x}_p$は一次独立です。

\paragraph{対称行列の固有値は実数}
一般に固有値や固有ベクトルを考える時は，$n$次方程式の根ですから複素数まで広げて考える必要があるのですが，幸いにして対称行列の場合はすべて実数で得られることがわかっています。

行列$\bm{A}$が実対称行列であるとし，この行列の固有値$\lambda$に対する固有ベクトルを$\bm{x}$としましょう。ここで$\lambda$や$\bm{x}$は複素数だとします。

\begin{quotation}
	複素数については本書の範疇を超えますが，少しだけその性質を紹介しておきます。$a,b$を実数，$i=\sqrt{-1}$として$a+bi$のことを複素数と言います。複素数$\alpha = a+bi$に対して$\overline{\alpha}=a-bi$を共役複素数といいます。その性質として，
	\begin{itemize}
		\item $\alpha$が実数\Leftrightarrow $\alpha$ = $\overline{\alpha}$
		\item $\overline{\alpha \pm \beta} = \overline{\alpha}\pm \overline{\beta}$
		\item $\overline{\alpha\beta} = \overline{\alpha}\overline{\beta}$
		\item $\alpha\overline{\alpha} = a^2+b^2$
	\end{itemize}
	などがあります。
\end{quotation}
$\bm{A}$は対称行列ですから，$\bm{A}=\bm{A}'$ですし，$\bm{A} = \overline{\bm{A}}$ でもあります。
さてこの固有値と固有ベクトルを使って，次のような計算を考えてみましょう。
\[
	\begin{array}{lllr}
		\lambda \overline{\bm{x}'}\bm{x} & = & \overline{\bm{x}'}\lambda \bm{x}           & \mbox{スカラーなので}  \\
		                                 &   &                                            &                 \\
		                                 & = & \overline{\bm{x}'}\bm{Ax}                  & \mbox{固有値の定義}   \\ & & &\\
		                                 & = & \overline{\bm{x}'}\overline{\bm{A}}\bm{x}  & \mbox{行列は実数なので} \\ & & &\\
		                                 & = & \overline{\bm{x}'}\overline{\bm{A}'}\bm{x} & \mbox{行列は対称なので} \\ & & &\\
		                                 & = & (\overline{\bm{Ax}})'\bm{x}                & \mbox{転置のルールから} \\ & & &\\
		                                 & = & (\overline{\lambda\bm{x}})'\bm{x}          & \mbox{固有値に戻す}   \\ & & &\\
		                                 & = & \overline{\lambda}\overline{\bm{x}}'\bm{x}                   \\
	\end{array}
\]
ということで，
\[ \lambda \overline{\bm{x}'}\bm{x} = \overline{\lambda}\overline{\bm{x}}'\bm{x} \]
という関係が得られましたが，これを移項するなどして
\[
	(\lambda -\bar{\lambda})\overline{\bm{x}}'\bm{x}=0.
\]
であることがわかります。ここで、$\bm{x} \neq \mbox{\boldmath $0$}$ですから，$\overline{\bm{x}}'\bm{x}\neq 0$より$\lambda = \bar{\lambda}$であることになります。ある値$\lambda$が複素数であるとしても，その共役複素数$\bar{\lambda}$と等しいことから、$\lambda$が実数であることが示ました。
\begin{reidai}{}{}
	対称行列$\bm{A} = \begin{pmatrix}1 & 2 \\ 2 & 3 \end{pmatrix}$の固有値と固有ベクトルを求めましょう。固有値が実数になることを確認してください。
\end{reidai}

\paragraph{対称行列の固有ベクトルは直交する}
さてこんどは同じく対称行列の，固有ベクトルに関する性質です。
対称行列$\bm{A}$の異なる固有値，$\lambda_1,\lambda_2$に対して，それぞれ固有ベクトル$\bm{x},\bm{y}$があるとします。
\[ \bm{Ax} = \lambda_1\bm{x},\bm{Ay} = \lambda_2\bm{y}\]
です。先ほどの式展開とほとんど同じなのですが，
\[
	\begin{array}{lllr}
		\lambda\bm{x}'\bm{y} & = & (\lambda \bm{x})']\bm{y} & \mbox{スカラーなので} \\
		                     &   &                          &                \\
		                     & = & (\bm{Ax})'\bm{y}         & \mbox{固有値の定義}  \\
		                     & = & \bm{x}'\bm{A}'\bm{y}     & \mbox{転置の計算}   \\
		                     & = & \bm{x}'\bm{Ay}           & \mbox{対称行列なので} \\
		                     & = & \bm{x}' \lambda_2\bm{y}  &                \\
		                     & = & \lambda_2 \bm{x}'\bm{y}
	\end{array}
\]
となり，
\[ (\lambda_1 - \lambda_2)\bm{x}'\bm{y} =0 \]
が得られます。$\lambda_1 \neq \lambda_2$なので，$\bm{x}'\bm{y} = 0$，内積が0なので直交していることがわかります。固有ベクトルを求めることは，(正規)直交基底を求めることでもあるのです。

\begin{reidai}{}{}
	さきほどの対称行列$\bm{A} = \begin{pmatrix}2 & 2 \\ 2 & -1 \end{pmatrix}$の固有ベクトルが直交していることを確認してください。
\end{reidai}
\paragraph{対角化とスペクトル分解}
固有値と固有ベクトルがわかると，行列をおもしろい形に変えることができます。

$\bm{A}$を$n$次正方対称行列，つまり$n \times n$サイズの対称行列だとすると，$n$個の固有値が求められます。これを$\lambda_1,\lambda_2,\cdots,\lambda_n$として，対応する固有ベクトルを$\bm{x}_1,\bm{x}_2,\cdots,\bm{x}_n$とします。ここで各固有ベクトルのノルムが1であるとしましょう。行列と固有値・固有ベクトルの関係から，
\[ \bm{Ax}_i = \lambda_i \bm{x}_i\]
となりますが，このベクトルを並べた行列$\bm{X}=\begin{pmatrix}\bm{x}_1&\bm{x}_2& \cdots & \bm{x}_n\end{pmatrix}$を考えると，
\[ \bm{AX} = \bm{X\Lambda} \]
と書くことができます。ここで$\bm{\Lambda}$は
\[ \bm{\Lambda} =
	\begin{pmatrix} \lambda_1 & 0         & \cdots & 0         \\
                0         & \lambda_2 & \cdots & 0         \\
                \vdots    & \vdots    & \ddots & \vdots    \\
                0         & 0         & \cdots & \lambda_n\end{pmatrix}
\]
のような行列です。

対称行列の異なる固有値に対する固有ベクトルは直交するので，$\bm{X}$は直交行列です。さてこの関係から，
\[ \bm{AXX}' = \bm{X\Lambda X}'\]
を考えれば，$\bm{XX}'=\bm{I}$ですから，
\[ \bm{A} = \bm{X\Lambda X}'\]
となることがわかります。正方行列は，自身の固有ベクトルからなる行列$\bm{X}$と，固有値からなる行列$\bm{\Lambda}$と，固有ベクトルの転置行列$\bm{X}'$の積に分解できることがわかりました。この分解のことを\keytermL{スペクトル分解}{すぺくとるぶんかい}とよぶことがありますが，この話はまた後で出てきますので，どうぞお楽しみに。

ついでに，さらにこの行列を固有ベクトルからなる行列で挟んで，
\[ \bm{X}'\bm{AX} = \bm{X}'\bm{X\Lambda X}'\bm{X} = \bm{\Lambda}\]
とすることもできます。$\bm{\Lambda}$は固有値が対角に入った対角行列ですから，このように変形することを\keytermL{対角化}{たいかくか}といいます。

この操作の何が楽しいかというと，行列の累乗計算がずいぶんと楽になるところです。$\bm{A}^n$を考えると，
\[ \begin{array}{lll}
		\bm{A}^n & = & \bm{AAA}\cdots\bm{A}                                                   \\
		         & = & \bm{X\Lambda X}'\bm{X\Lambda X}'\bm{X\Lambda X}'\cdots\bm{X\Lambda X}' \\
		         & = & \bm{X} \bm{\Lambda\Lambda\cdots\Lambda} \bm{X}'                        \\
		         & = & \bm{X} \bm{\Lambda}^n \bm{X}'                                          \\
		         &   &                                                                        \\
		         & = & \bm{X}
		\begin{pmatrix} \lambda_1^n & 0           & \cdots & 0           \\
                0           & \lambda_2^n & \cdots & 0           \\
                \vdots      & \vdots      & \ddots & \vdots      \\
                0           & 0           & \cdots & \lambda_n^n\end{pmatrix}\bm{X}'
	\end{array}
\]
となるからです。行列$\bm{A}$が座標の変換であったり，何らかの作用を表現した行列であるとすると，それを作用させ続けた場合を考える時に累乗計算になりますが，その行列計算が一気に求められることになります。「何らかの作用」をハミルトン演算子と見るか，ページランクとみるか\footnote{Googleの検索エンジンが採用しているページの評価法。Googleの検索システムは，ページ同士のリンク関係を行列で表現し，スペクトル分解して重要度の高いページを重みづけて表現する，というアルゴリズムをとっており，「検索したいものに効率よく辿り着ける」ことで人気を博したからです。}，隣接行列\footnote{数理社会学などでは対人的なネットワークの関係を行列で表現し，このように呼びます。ネットワーク科学一般でも用いられてる用語と言えるでしょう。}と見るかによって見える景色は違ってくるかもしれませんが，数学的には固有値問題といういみで同じなのです。興味がある人は\citet{Naganuma20110920}の電子版付録などもみてみるといいでしょう。

\paragraph{総和はトレースに一致する}
\label{sumIsTrace}
さて今度は，$n$個の固有値を$\lambda_1,\lambda_2,\cdots,\lambda_n$とすると
\[ \sum_i^n \lambda_i = tr(\bm{A}) \]
ということを示します。これは対角化できることを考えると，
\[ tr(\bm{A}) = tr(\bm{X\Lambda X}') \]
となります。トレースは行列の積に対して順序の入れ替えを許しますから($\to$Pp.\pageref{defineOfTrace})，
\[ =tr(\bm{\Lambda XX}') = tr(\bm{\Lambda}) = \sum_i^n \lambda_i \]
です。

\begin{reidai}{}{}
	さきほどの対称行列$\bm{A} = \begin{pmatrix}2 & 2 \\ 2 & -1 \end{pmatrix}$の固有値の和がトレースと一致することを確認してください。
\end{reidai}

\section{固有値と固有ベクトルの幾何学的意味}

固有値，固有ベクトルについて，今度は違う側面から見直してみましょう。

ある正方行列から固有値$\lambda$と固有ベクトル$\bm{a}$が得られたとします。このベクトル$\bm{a}$のすべての要素を定数$c$倍したベクトル$\bm{b}=c\bm{a}$を考えると，これもやはり同じ関係が成り立ちます。
\begin{equation}
	\bm{A}\bm{b}=  \lambda c \bm{a} =\lambda \bm{b}
\end{equation}
先ほどの計算でも明らかになりましたが，固有ベクトルの値は絶対的なものではなく，要素間の相対的大きさを反映しているに過ぎないのでしたね。

さてこれを幾何学的に，図形として考えてみましょう。要素が2つのベクトルは，2次元座標に表現できます。ベクトル$\bm{x}=(x,y)$という座標を表しているというわけです。固有ベクトルも要素が2つであれば，座標で表現できます。先ほどの，要素を$c$倍しても固有ベクトルとしての性質は変わらない，という話は，「固有ベクトルは大きさに意味はなく，方向を表したもの」ということになります。では何の方向を指し示しているのでしょうか。

ここで行列の空間イメージで出てきた，「行列は座標変換装置」ということを思い出してください。
今回の例では，行列$\bm{A}$には次のような性質があります。
\begin{equation}
	\begin{pmatrix}
		2 & 0 \\
		0 & 3
	\end{pmatrix}
	\begin{pmatrix}
		1 \\
		0
	\end{pmatrix}
	= 2 \begin{pmatrix} 1 \\ 0 \end{pmatrix}
\end{equation}

\begin{equation}
	\begin{pmatrix} 2 & 0\\ 0 & 3 \end{pmatrix} \begin{pmatrix} 0 \\ 1 \end{pmatrix}
	= 3 \begin{pmatrix} 0 \\ 1 \end{pmatrix}
\end{equation}

そう，お気づきのように，これは固有値・固有ベクトルです。この行列の固有値・固有ベクトルはそれぞれ$\lambda_1=2,\bm{x}_1=(1,0)$，$\lambda_2=3,\bm{x}_2=(0,1)$であることがわかります。この固有値，固有ベクトルの組み合わせをじっとみていると，おもしろい特徴がわかって来ます。

今回の行列から得られた2つの固有ベクトル，$(1,0)$と$(0,1)$は，2次元平面の正規直交基底であることにお気づきのことと思います。つまり，正方行列$\bm{A}$から得られる固有ベクトルは，その正方行列がもつ基底だといえるのです。

では固有値はどうでしょうか？ 今回はx座標を2倍，y座標を3倍に引き延ばす変換をしたわけですが，この座標の歪み(重み)が固有値に対応していますね。つまり，固有値と固有ベクトルは新しい座標に変換する，その変換先の空間的性質を表していることになります。元の座標空間は$(1,0),(0,1)$で作られる空間ですが，変換先の空間は$(2,0),(0,3)$で作られている空間，ということになります。

つまり，正方行列は空間を変換するものであり，正方行列がもっている固有ベクトルを基底とした空間を通して変換されるということになります。


\vspace{2cm}
固有値と固有ベクトルの関係は，次にお話しする因子分析の話に直結します。多変量解析の文脈で行列がどのように応用されるのか，みていきましょう。
\end{document}